\chapter{Fine-tuning des LLM}
\label{ch:finetuning}

\begin{quote}
\emph{<< Le pr\'e-entra\^inement donne au mod\`ele de la connaissance. Le fine-tuning lui donne un m\'etier. >>}
\end{quote}

% --- hook ---
Avant 2021, fine-tuner un modèle de 7 milliards de paramètres demandait plusieurs GPU A100 et un week-end entier. Edward Hu, alors stagiaire chez Microsoft, propose une idée qui paraît trop simple pour fonctionner : ne pas modifier les poids du modèle, mais leur ajouter une petite matrice de rang faible. C'est LoRA. Deux ans plus tard, Dettmers et son équipe étendent l'idée en chargeant le modèle de base en 4~bits --- QLoRA --- et soudain, un Llama de 7B se fine-tune sur un GPU Colab gratuit. Ce chapitre raconte comment la barrière d'entrée du fine-tuning a chuté d'un facteur 50 en deux ans, et comment vous, sur un T4 emprunté, pouvez en bénéficier.
\section{Pourquoi fine-tuner ?}

Le prompt engineering a ses limites. Quand vous avez besoin qu'un mod\`ele :
\begin{itemize}
  \item Suive constamment un format de sortie sp\'ecifique,
  \item Manipule un vocabulaire de domaine (juridique, m\'edical, scientifique),
  \item Atteigne une pr\'ecision maximale sur une t\^ache \'etroite,
  \item R\'eduise la latence avec un mod\`ele plus petit et sp\'ecialis\'e,
\end{itemize}
le fine-tuning est la r\'eponse.

% ============================================================
\section{Fine-tuning complet vs m\'ethodes \`a param\`etres efficaces}

\textbf{Fine-tuning complet} met \`a jour tous les param\`etres du mod\`ele. Pour un mod\`ele 7B, cela demande :
\begin{itemize}
  \item 28 Go pour les poids (FP32)
  \item 28 Go pour les \'etats de l'optimiseur (Adam)
  \item 28 Go pour les gradients
  \item Total : $\sim$84 Go de VRAM --- plusieurs GPU A100.
\end{itemize}

\textbf{Fine-tuning \`a param\`etres efficaces (PEFT)} gel la plupart des param\`etres et n'entra\^ine qu'un petit nombre de param\`etres suppl\'ementaires.

% ============================================================
\section{LoRA : Low-Rank Adaptation}

LoRA d\'ecompose les mises \`a jour de poids en matrices de bas rang. Au lieu de mettre \`a jour $W \in \mathbb{R}^{d \times d}$, on apprend $\Delta W = BA$ avec $B \in \mathbb{R}^{d \times r}$ et $A \in \mathbb{R}^{r \times d}$, et un rang $r \ll d$.

\[
W' = W + \alpha \cdot BA
\]

Avantages : seulement $2dr$ param\`etres au lieu de $d^2$. Pour $d=4096, r=16$ : on passe de 16,7M \`a 131K param\`etres par couche.

\begin{minted}{python}
from peft import LoraConfig, get_peft_model, TaskType
from transformers import AutoModelForCausalLM

model = AutoModelForCausalLM.from_pretrained(
    "TinyLlama/TinyLlama-1.1B-Chat-v1.0",
    torch_dtype="auto",
    device_map="auto",
)

lora_config = LoraConfig(
    task_type=TaskType.CAUSAL_LM,
    r=16,                        # rank
    lora_alpha=32,               # scaling factor
    lora_dropout=0.05,
    target_modules=["q_proj", "v_proj", "k_proj", "o_proj"],
)

peft_model = get_peft_model(model, lora_config)
peft_model.print_trainable_parameters()
# trainable params: ~4M / total: ~1.1B = 0.36%
\end{minted}

% ============================================================
\section{QLoRA : LoRA quantifi\'e}

QLoRA charge le mod\`ele de base en pr\'ecision 4 bits (quantification NF4), puis entra\^ine les adapteurs LoRA en FP16. Cela r\'eduit l'usage de VRAM d'un facteur $\sim$4 :

\begin{minted}{python}
from transformers import BitsAndBytesConfig
import torch

bnb_config = BitsAndBytesConfig(
    load_in_4bit=True,
    bnb_4bit_quant_type="nf4",
    bnb_4bit_compute_dtype=torch.float16,
    bnb_4bit_use_double_quant=True,
)

model = AutoModelForCausalLM.from_pretrained(
    "TinyLlama/TinyLlama-1.1B-Chat-v1.0",
    quantization_config=bnb_config,
    device_map="auto",
)

# Appliquer LoRA par-dessus le mod\`ele quantifi\'e
peft_model = get_peft_model(model, lora_config)
peft_model.print_trainable_parameters()
\end{minted}

\begin{aitip}
QLoRA rend possible le fine-tuning d'un mod\`ele 7B sur un seul GPU Colab T4 gratuit (16 Go de VRAM). Cela d\'emocratise le fine-tuning --- pas besoin de mat\'eriel co\^uteux.
\end{aitip}

% ============================================================
\section{Jeux de donn\'ees pour fine-tuning}

\begin{longtable}{p{3cm}p{3cm}p{6.5cm}}
\toprule
\textbf{Dataset} & \textbf{Taille} & \textbf{Description} \\
\midrule
Alpaca & 52K & Donn\'ees instruction-following de Stanford \\
Dolly & 15K & Instructions r\'edig\'ees par des employ\'es Databricks \\
OpenAssistant & 160K & Conversations multi-tours \\
UltraChat & 1.5M & Dialogues multi-tours synth\'etiques \\
\bottomrule
\end{longtable}

\begin{minted}{python}
from datasets import load_dataset

# Charger le dataset Dolly
dataset = load_dataset("databricks/databricks-dolly-15k", split="train")
print(f"Dataset size: {len(dataset)}")
print(dataset[0])

# Formater pour l'instruction tuning
def format_instruction(example):
    if example["context"]:
        text = (f"### Instruction:\n{example['instruction']}\n\n"
                f"### Context:\n{example['context']}\n\n"
                f"### Response:\n{example['response']}")
    else:
        text = (f"### Instruction:\n{example['instruction']}\n\n"
                f"### Response:\n{example['response']}")
    return {"text": text}

formatted = dataset.map(format_instruction)
print(formatted[0]["text"][:300])
\end{minted}

% ============================================================
\section{Fine-tuning avec HuggingFace Trainer}

\begin{minted}{python}
from transformers import (
    AutoModelForCausalLM, AutoTokenizer,
    TrainingArguments, Trainer, DataCollatorForLanguageModeling,
    BitsAndBytesConfig,
)
from peft import LoraConfig, get_peft_model, TaskType
from datasets import load_dataset
import torch

# 1. Charger le mod\`ele en 4 bits
bnb_config = BitsAndBytesConfig(
    load_in_4bit=True,
    bnb_4bit_quant_type="nf4",
    bnb_4bit_compute_dtype=torch.float16,
)

model_name = "TinyLlama/TinyLlama-1.1B-Chat-v1.0"
tokenizer = AutoTokenizer.from_pretrained(model_name)
tokenizer.pad_token = tokenizer.eos_token

model = AutoModelForCausalLM.from_pretrained(
    model_name, quantization_config=bnb_config, device_map="auto"
)

# 2. Appliquer LoRA
lora_config = LoraConfig(
    task_type=TaskType.CAUSAL_LM, r=16, lora_alpha=32,
    lora_dropout=0.05,
    target_modules=["q_proj", "v_proj", "k_proj", "o_proj"],
)
model = get_peft_model(model, lora_config)

# 3. Pr\'eparer le dataset
dataset = load_dataset("databricks/databricks-dolly-15k", split="train[:1000]")

def tokenize(example):
    text = f"### Instruction:\n{example['instruction']}\n### Response:\n{example['response']}"
    return tokenizer(text, truncation=True, max_length=512, padding="max_length")

tokenized = dataset.map(tokenize, remove_columns=dataset.column_names)

# 4. Param\`etres d'entra\^inement
training_args = TrainingArguments(
    output_dir="./tinyllama-dolly-lora",
    num_train_epochs=3,
    per_device_train_batch_size=4,
    gradient_accumulation_steps=4,
    learning_rate=2e-4,
    fp16=True,
    logging_steps=10,
    save_strategy="epoch",
    warmup_ratio=0.03,
    lr_scheduler_type="cosine",
)

# 5. Entra\^iner
trainer = Trainer(
    model=model,
    args=training_args,
    train_dataset=tokenized,
    data_collator=DataCollatorForLanguageModeling(tokenizer, mlm=False),
)
trainer.train()

# 6. Sauvegarder l'adapteur LoRA
model.save_pretrained("./tinyllama-dolly-lora")
\end{minted}

% ============================================================
\section{\'Evaluation apr\`es fine-tuning}

\begin{minted}{python}
# Charger le mod\`ele fine-tun\'e
from peft import PeftModel

base_model = AutoModelForCausalLM.from_pretrained(
    model_name, quantization_config=bnb_config, device_map="auto"
)
finetuned = PeftModel.from_pretrained(base_model, "./tinyllama-dolly-lora")

# Comparer base vs fine-tun\'e
prompt = "### Instruction:\nExplain what a neural network is.\n### Response:\n"
inputs = tokenizer(prompt, return_tensors="pt").to(model.device)

# Mod\`ele de base
base_out = base_model.generate(**inputs, max_new_tokens=100, temperature=0.7)
print("=== Base model ===")
print(tokenizer.decode(base_out[0], skip_special_tokens=True))

# Mod\`ele fine-tun\'e
ft_out = finetuned.generate(**inputs, max_new_tokens=100, temperature=0.7)
print("\n=== Fine-tuned model ===")
print(tokenizer.decode(ft_out[0], skip_special_tokens=True))
\end{minted}

\begin{warningbox}
Le fine-tuning sur de petits datasets risque l'\emph{oubli catastrophique} (catastrophic forgetting) : le mod\`ele perd ses capacit\'es g\'en\'erales. Utilisez un ensemble de validation r\'eserv\'e et surveillez \`a la fois la performance sp\'ecifique \`a la t\^ache et la performance g\'en\'erale.
\end{warningbox}

\begin{exercise}
\begin{enumerate}
  \item Fine-tunez TinyLlama sur 500 exemples du dataset Dolly avec QLoRA. Reportez la perte d'entra\^inement \`a chaque \'epoque.
  \item Exp\'erimentez avec les valeurs de rang LoRA : $r \in \{4, 8, 16, 32\}$. Lequel donne le meilleur compromis entre qualit\'e et temps d'entra\^inement ?
  \item Fine-tunez sur un dataset personnalis\'e : cr\'eez 100 paires instruction-r\'eponse sur un sujet qui vous int\'eresse (cuisine, histoire, ou un domaine sp\'ecifique). \'Evaluez le mod\`ele avant et apr\`es.
  \item Comparez l'usage de VRAM entre fine-tuning complet et QLoRA sur TinyLlama. Utilisez \texttt{torch.cuda.max\_memory\_allocated()}.
\end{enumerate}
\end{exercise}

\begin{ethicsbox}
Le fine-tuning peut int\'egrer des biais des donn\'ees d'entra\^inement. Si votre dataset contient des st\'er\'eotypes, le mod\`ele les apprendra. Auditez toujours vos donn\'ees pour rep\'erer les biais, testez le mod\`ele fine-tun\'e sur des entr\'ees diversifi\'ees, et documentez les sources de donn\'ees et limites connues.
\end{ethicsbox}

% ============================================================
\section{R\'esum\'e du chapitre}

\begin{itemize}
  \item Le fine-tuning adapte un mod\`ele pr\'e-entra\^in\'e \`a une t\^ache ou un domaine sp\'ecifique.
  \item Le fine-tuning complet met \`a jour tous les param\`etres et demande beaucoup de VRAM.
  \item LoRA entra\^ine des matrices d'adaptation de bas rang, r\'eduisant les param\`etres entra\^inables de plus de 99\%.
  \item QLoRA combine quantification 4 bits et LoRA, permettant le fine-tuning sur GPU Colab gratuit.
  \item Les datasets d'instruction-tuning (Alpaca, Dolly) apprennent aux mod\`eles \`a suivre des instructions.
  \item HuggingFace Trainer + PEFT fournit un pipeline complet de fine-tuning en $\sim$30 lignes.
  \item \'Evaluez toujours sur des donn\'ees r\'eserv\'ees et surveillez l'oubli catastrophique.
\end{itemize}
